\begin{abstract}
A widespread practice in software development is to tailor coding agents to repositories using \agentsmds{}, such as \texttt{AGENTS.md}, by either manually or automatically generating them. Although this practice is strongly encouraged by agent developers, there is currently no rigorous investigation into whether such \agentsmds{} are actually effective for real-world tasks. In this work, we study this question and evaluate coding agents' task completion performance in two complementary settings: established SWE-bench tasks from popular repositories, with LLM-generated \agentsmds{} following agent-developer recommendations, and a novel collection of issues from repositories containing developer-committed \agentsmds{}.

Across multiple coding agents and LLMs, we find that \agentsmds{} tend to \emph{reduce} task success rates compared to providing no repository context, while also \emph{increasing inference cost} by over 20\%. Behaviorally, both LLM-generated and developer-provided context files encourage broader exploration (e.g., more thorough testing and file traversal), and coding agents tend to respect their instructions. Ultimately, we conclude that unnecessary requirements from \agentsmds{} make tasks harder, and human-written context files should describe only minimal requirements.
\end{abstract}
